\documentclass[11pt,a4paper]{article}
\usepackage[T1]{fontenc}
\usepackage{lmodern,amsmath,amssymb,amsthm,mathtools,booktabs,microtype,xcolor}
\usepackage[margin=26mm]{geometry}
\usepackage[colorlinks=true,linkcolor=blue!45!black,citecolor=blue!45!black,urlcolor=blue!45!black]{hyperref}
\newtheorem{theorem}{Theorem}[section]
\newtheorem{lemma}[theorem]{Lemma}
\newtheorem{proposition}[theorem]{Proposition}
\newtheorem{corollary}[theorem]{Corollary}
\theoremstyle{remark}\newtheorem{remark}[theorem]{Remark}
\newcommand{\Q}{\mathbb Q}\newcommand{\Z}{\mathbb Z}
\newcommand{\HH}{\mathcal H}
\setlength{\emergencystretch}{2em}
\title{A quartic in 370 variables with specified integer coefficients}
\author{}\date{8 November 2026}
\begin{document}\maketitle
\begin{abstract}
We specify an integer quartic in 370 total variables whose rational
input projection is the integers, under the arithmetic premises of
the specified elliptic source construction. The coefficient table
contains 5837 nonzero monomials. The construction uses three height
detectors applied to rational linear forms in a sum of four
elliptic abscissas. An effective bound for the height lost in that
sum supplies an explicit numerical parameter. All four point charts
allow arbitrary nonzero integer indices. The resulting quartic has
content one and a largest coefficient with 1791357 binary digits.
We also specify a family with the same defining property for every
integer parameter above an explicit threshold.
\end{abstract}

\section{The polynomial and its arithmetic premises}

For a polynomial $F\in\Z[X,Y_1,\ldots,Y_m]$, its rational input
projection is the set
\[
 \{x\in\Q:\exists (y_1,\ldots,y_m)\in\Q^m,\ F(x,\boldsymbol y)=0\}.
\]
All variable counts include $X$. Degree is total degree. The height
$H(F)$ of an integer polynomial is its largest absolute coefficient.
For a rational or algebraic number, $h$ denotes absolute logarithmic
Weil height.

Put $\ell=80989$. The exact tables
\nolinkurl{quadratic-family370.json} and
\nolinkurl{polynomial-family370.json} specify
\[
 A_N=\ell^{2N+1},\qquad
 P_N=\sum_{j=1}^{342}Q_{N,j}^2
       \in\Z[X,Y_1,\ldots,Y_{369}],
\]
where every $Q_{N,j}$ is quadratic. The table records the powers
of $A_N$ separately to keep large integer coefficients readable;
$A_N$ is a coefficient, not a variable.

\begin{theorem}\label{thm:main}
Assume the arithmetic premises stated below. For every integer
$N\ge27464$ and every $x\in\Q$,
\[
 x\in\Z\quad\Longleftrightarrow\quad
 \exists\boldsymbol y\in\Q^{369},\ P_N(x,\boldsymbol y)=0.
\]
Every such polynomial has degree exactly four, content one,
and 5837 nonzero monomials. Its coefficient height satisfies
\[
                   H(P_N)<2^{1520}\,80989^{4N+2}.
\]
In particular, the single polynomial $P_{27464}$ is fully specified
by substituting $A=80989^{54929}$ in the table. Its largest
coefficient has exactly 1791357 binary digits.
\end{theorem}

We retain the arithmetic contract of the explicit source in
\cite[Section 2]{Quartic392}. In particular, with
\[
 K=\Q(\rho),\quad\rho^2=2,\quad d=3-\rho,\quad
 E:y^2=x^3-d^2x,
\]
\[
 R_0=(1+\rho,2\rho),\qquad
 P=2R_0=\left(\frac{57-28\rho}{8},
                \frac{-660+417\rho}{32}\right),
\]
the group and index-height premises are
\begin{equation}\label{eq:group}
 2E(K)=\Z P,\qquad h(x([n]P))\ge n^2-6\quad(n\in\Z\setminus\{0\}).
\end{equation}
The contact-height premise, pole obstruction and simultaneous pole
realization, and simultaneous prime-pattern premise are unchanged.
The quantitative weighted selection premise used in the source's
real-separation argument is also retained. These inputs are explicitly
enumerated in \cite[Section 2]{Quartic392}; coefficient verification
does not establish them.

For reference, the contact-height constant is
\begin{equation}\label{eq:H}
 \HH=2^{110}210^{1800}
          \left(\prod_{p\le76176}p\right)^{19}+128
       <2^{2093448}.
\end{equation}
The effective abscissa-sum theorem used below is proved separately
in \cite{EffectiveSums}. It is not an additional conjectural premise,
and it does not require an ineffective height constant.

\section{Point charts and the three height scales}

The four root charts retain the quadratic equations over $K$
\begin{equation}\label{eq:chart}
                         x=h^2,\qquad j^2=x^2-d^2.
\end{equation}
The corresponding elliptic point is $Q=(x,hj)$.

\begin{lemma}\label{lem:indices}
Under \eqref{eq:group}, each chart \eqref{eq:chart} gives a point
$Q=[m]P$ with $m$ a nonzero integer. Conversely every nonzero
integer multiple of $P$ admits such a chart. No parity restriction
on $m$ is needed.
\end{lemma}
\begin{proof}
The field $K$ is real, so $h=0$ would imply $j^2=-d^2$, an
impossibility. Thus the chart gives a finite nonzero point of $E$.
By \eqref{eq:group}, write $Q=mR_0+T$ for $T\in E[2](K)$.
The standard two-isogeny squareclass homomorphism has value
$[x(Q)]$ away from the origin and $(0,0)$, value $1$ at the origin,
and value $[-d^2]$ at $(0,0)$. Its values on $E[2](K)$ are
$1,-d^2,d,-d$. These four classes have uniform signs at the
two real embeddings, whereas $[x(R_0)]=[1+\rho]$ has mixed
signs. Since $x(Q)$ is a square, $m$ is even. The torsion class
must then be a square. Neither negative class is a square, and
$d$ is not a square because its rational norm is $7$. Hence
$T=O$ and $Q=[m/2]P$; finiteness makes this index nonzero.

Conversely every such point is a double. The duplication formula
\[
 x(2V)=\left(\frac{x(V)^2+d^2}{2y(V)}\right)^2
\]
gives $x=h^2$, and putting $j=y/h$ gives the second chart equation.
The relevant numerator and denominator cannot vanish for a
nontorsion point. The squareclass argument and all exceptional
point exclusions are also detailed in \cite[Section 2]{Quartic317}.
\end{proof}

Write the four root points as
\[
 Q_1=[n]P,\qquad Q_2=[a]P,\qquad
 Q_3=[t]P,\qquad Q_4=[b]P,
 \qquad B=\max(|n|,|a|,|t|,|b|).
\]
Define
\[
 S=x(Q_1)+x(Q_2)+x(Q_3)+x(Q_4)=S_0+S_1\rho,
 \qquad S_0,S_1\in\Q.
\]
Every abscissa is positive at both real embeddings. In fact its
curve equation and its nonzero ordinate give $x(Q_i)>d$ at each
embedding. Consequently $S_0>12$ and $S_0>\sqrt2|S_1|$.
In particular, both $S_0$ and $S_0+S_1$ are positive.

\begin{lemma}\label{lem:height}
For every tuple supplied by the root charts,
\[
                  B^2\le2h(S_0)+2h(S_0+S_1)+82.
\]
\end{lemma}
\begin{proof}
The effective sum-height theorem \cite{EffectiveSums} states
\[
 h(S)\ge\max_i h(x(Q_i))-24\log B-\log864.
\]
Its proof compares the integral local principal parts
$x([n]P)=(nL_v)^{-2}+O(1)$ at each finite place.
For a fixed target index, only subsets containing that index
can cause loss; their positive integer numerators have total
size bounded by $B^{24}\cdot864$. Positivity prevents loss
at the real places. Applying \eqref{eq:group} and the elementary
logarithm bound gives $B^2\le2h(S)+80$.
Write $U=S_0$ and $V=S_0+S_1$, so
\[
                         S=(1-\rho)U+\rho V.
\]
Both coefficients are algebraic integers. Hence the finite-place
part of $h(S)$ is at most that of $h(1,U,V)$. At the two real
embeddings, the sums of their absolute values are $2\sqrt2-1$
and $2\sqrt2+1$, whose product is $7$. Therefore
\[
 h(S)\le h(1,U,V)+\tfrac12\log7
       \le h(U)+h(V)+\tfrac12\log7.
\]
Since $\log7<2$, the assertion follows.
\end{proof}

Use the three positive rational square scales
\begin{equation}\label{eq:scales}
                      1,\qquad S_0^2,\qquad(S_0+S_1)^2.
\end{equation}
All point tests, logarithm approximations, comparison tests and
pole equations in the three detectors retain their exact source
forms. Their orders are $1,10,40$. The local nonvanishing argument
of \cite{Quartic389} applies to arbitrary nonzero square scales,
including these three scales; it covers the degenerate input
$t=1$ and empty contact radicals. Thus no deleted inverse
witness is reintroduced.

\begin{lemma}\label{lem:divheight}
Suppose $J$ is a nonzero integer and the $k$-th powers of all
three contact radicals divide $J$. Then
\[
                  B^2\le5\HH|J|^{19/k}.
\]
\end{lemma}
\begin{proof}
Put $T=(\HH+\log3)|J|^{19/k}$. The retained contact-height
theorem bounds each of
\[
 h(\lambda),\quad h(S_0^2\lambda),\quad
 h((S_0+S_1)^2\lambda)
\]
by $T$; its fractional linear contact change accounts for
$\log3$. Since the detector domain excludes $\lambda=0$,
\[
                    h(S_0)\le T,\qquad h(S_0+S_1)\le T.
\]
Lemma~\ref{lem:height} now gives
\[
                            B^2\le4T+82.
\]
Here $|J|\ge1$ and the explicit constant \eqref{eq:H} is
greater than $92$. Therefore $4T+82\le5\HH|J|^{19/k}$.
\end{proof}

\section{The explicit scalar parameter and soundness}

The input label and scalar polynomial are
\[
 \beta=51X^2+2504,\qquad
 r=\beta(\beta-1)(\beta-2401)(32\beta+343)(2\beta-1),
\]
\[
 R=\frac{A_Nr}{1+\ell^2r},\qquad c_R=R+9792.
\]
The retained scalar block, with the indicated shift, is precisely
the parity-free scalar block of \cite[Sections 4--5]{Quartic317}.
Its integral group identity is
$[m](2G)=[4(a-1)]H$; it does not require $a-1$ to be even.
The existing exceptional-point exclusions and denominator guards
are retained.

For clarity, its conclusions before any height comparison are
\begin{equation}\label{eq:scalar}
 c:=|a-1|\ge\ell^{L_N},\qquad
 L_N=\left\lfloor\frac{187(2N-1)}8\right\rfloor+1,
 \qquad h(\beta)\le\HH c^{19/24}.
\end{equation}
The proof uses the same pole obstruction, the same formal
comparison and the same five contact values as the source.
The fixed prime avoids these five squareclasses: for each
$\theta\in\{0,1,2401,-343/32,1/2\}$, the rational number
$(\theta-2504)/51$ is a nonsquare modulo $80989$.
The exact Euler tests are repeated by the checker.

At $N=27464$ one has $L_N=1283919$. The integer certificates
\[
 80989^{4096}>2^{66787},\qquad
 66787\cdot1283919-4096\cdot20934550=1181453>0
\]
and \eqref{eq:H} imply, for every $N\ge27464$,
\begin{equation}\label{eq:numerical}
 c>2^{20934550}>(80\HH)^{10},\qquad
 |a|>(40\HH)^{10},\qquad
 \log(48\HH)<|a|/40,
 \qquad h(\beta)<|a|/40.
\end{equation}
These are the same elementary numerical implications as in
\cite[Section 4]{Quartic392}. For example,
$c^{5/24}>80\HH$ gives $h(\beta)<c/80\le |a|/40$,
and the remaining estimates follow directly from $c-1\le|a|$.

The unchanged comparison equations imply the following
divisibilities. Put
\[
 \Delta=t^2-a(a-2)tb+(a-2)^2b^2-1.
\]
The fortieth powers of all three contact radicals divide
$\Delta$. If $n-\beta\ne0$, their tenth powers divide the
nonzero reduced numerator $J$ of $n-\beta$. These are the
local comparison conclusions of the source; removing parity
certificates does not alter their proofs, as established in
\cite{Quartic317}.

\begin{proposition}\label{prop:soundness}
Every rational zero of $P_N$, for $N\ge27464$, has integer input.
\end{proposition}
\begin{proof}
Over $\Q$, the sum of squares vanishes precisely when every
quadratic row vanishes. Suppose first that $\Delta\ne0$.
For $B\ge1$,
\[
 |\Delta|\le2B^4+6B^3+5B^2+1\le14B^4.
\]
Lemma~\ref{lem:divheight} gives
\[
 B^2\le5\HH(14B^4)^{19/40}<20\HH B^{19/10},
\]
because $14^{19}<4^{40}$. Hence $B^{1/10}<20\HH$,
contradicting $B\ge|a|>(40\HH)^{10}$.
Thus $\Delta=0$.

The elementary Lucas descent proved in
\cite[Section 4]{Quartic392} applies to all nonzero integer
$b$ and $|a|>7$, without a parity condition. It gives
\[
                            \log B>|a|-3.
\]
If $n-\beta\ne0$, write $J$ for its reduced numerator.
The elementary rational-height bound gives
$|J|\le2B e^{h(\beta)}$. The order-ten divisibility and
Lemma~\ref{lem:divheight} would imply
\[
 B^2\le5\HH(2B e^{h(\beta)})^{19/10}
       <20\HH B^{19/10}e^{19h(\beta)/10},
\]
and therefore
\[
 \log B<10\log(20\HH)+19h(\beta)
       <\frac{29}{40}|a|.
\]
This contradicts the Lucas bound because $|a|>11$.
Consequently $n=\beta\in\Z$.

Finally write $X=u/v$ in lowest terms, $v>0$. The equality
$51X^2+2504\in\Z$ implies $v^2\mid51$. Since $51$ is
squarefree, $v=1$, proving $X\in\Z$.
\end{proof}

\section{Completeness and equal scales}

\begin{proposition}\label{prop:complete}
Under the stated arithmetic premises, every integer input
extends to a rational zero of $P_N$ for every $N\ge27464$.
\end{proposition}
\begin{proof}
For an integer input, $n=\beta=51X^2+2504$ is an integer.
Choose the remaining integer labels and elliptic witnesses
by the scalar and Lucas completeness construction in
\cite{Quartic317}. The root charts admit every nonzero
integer label by Lemma~\ref{lem:indices}.

The small-prime scalar checks can be stated explicitly.
For the integer $r$ above,
\[
 v_2(r)=\begin{cases}3,&4\mid X,\\2,&4\nmid X,\end{cases}
 \qquad v_3(r)=1,\qquad r\equiv15\pmod{17}.
\]
Since $9792=2^6\cdot3^2\cdot17$, the local correction
$(R+9792)/R$ is a square at $2,3,17$. These are exactly
the local conditions used in the retained scalar pole
realization. The fixed avoidance at $\ell$ was already
checked above.

After fixing the integer root labels, apply the retained
prime-pattern construction to the distinct scales in
\eqref{eq:scales}. They are positive rational squares, and
at least two remain since $S_0>12$ implies $S_0^2\ne1$.
Equal scales have identical mathematical detectors and
reuse the same witnesses. No equation requires the three
scales to be distinct; in particular $S_1=0$ creates no
exception.

The conjugate reduction-order congruences used for the
comparison points remain valid for arbitrary fixed integer
root labels. At the fixed small primes, choose compact
local neighborhoods avoiding the finitely many poles of
the rational functions, as in the source completeness proof.
Their valuations are then bounded below by constants
depending on the already fixed labels. The pole
preparation can increase the chosen local depths after these
constants have been fixed. This is the same local argument
used for even input labels in \cite{Quartic317}; it needs
no uniform bound as the input varies. The real approximation
and quantitative separation steps then give the simultaneous
height-pole witnesses for these three scales.

Fill in the remaining arithmetic-circuit coordinates by
their stated rational formulas. The retained denominator
conditions hold for these completeness witnesses. The two
new scale equations hold by their definitions in the next
section. Hence all 342 quadratic rows vanish, as required.
\end{proof}

\section{Literal specification and independent verification}

The exact source system is the 389-variable table of
\cite{Quartic389}, whose uncompressed SHA-256 digest is
\begin{center}\small
\nolinkurl{f7fb5bdf6a3d47275f1d78ed5e275954e2cc1f21fec8204d6a395bfc2f8b5394}.
\end{center}
Source identifiers are stable identifiers in that table;
they are not the consecutive coordinate numbers of the
new system. The complete list of all 370 coordinates is
included in both new tables.

For the source root coordinates, the actual abscissas are
\[
 \begin{array}{ll}
 x(Q_1)=x_{82}+2x_{83}\rho,&x(Q_2)=x_{84}+2x_{85}\rho,\\
 x(Q_3)=2x_{86}+2x_{87}\rho,&x(Q_4)=x_{88}+2x_{89}\rho.
 \end{array}
\]
Thus
\[
 S_0=x_{82}+x_{84}+2x_{86}+x_{88},\qquad
 S_1=2(x_{83}+x_{85}+x_{87}+x_{89}).
\]
The two scale coordinates satisfy
\[
                       S_0^2=4x_{343},\qquad
                       (S_0+S_1)^2=4x_{344}.
\]
For the third detector, the constructor uses a separate coordinate
copy of the entire second detector, with $x_{343}$ replaced by
$x_{344}$. Both nonconstant detectors therefore use the same
literal normalization: their scale is four times the stored
scale coordinate and their contact input is four times their
stored input coordinate. The semantic scales are exactly
\eqref{eq:scales}.

More explicitly, for either nonconstant scale $a$, put
$t=a\lambda$ and $\delta=t-1$, and let
\[
 \Theta=\{0,-1/2,-1,1/2399,-32/407\},\qquad
 F(t)=\prod_{\theta\in\Theta}(t-\theta).
\]
The independently expanded equations are
\[
 F(t)e=(a\delta)^5,\qquad \delta^3q_1=e,\qquad
 e\delta b=T^2.
\]
The retained toric gates compute
$q_{10}=e^{19}/\delta^3$ and $q_{40}=e^{79}/\delta^3$
on their nonzero domain. Thus every power, field product and
pole input has the source normalization. At consumed contact
primes the local nonvanishing argument justifies that domain.

The constructor removes the six product coordinates with
source identifiers $337,\ldots,342$ and their six equations.
It also removes the twelve odd-index coordinates and their
twelve sole-use equations, and the obsolete scale coordinate
$x_{345}$ and its defining equation. The old input is replaced by
$(X-1)/2$, its square coordinate by $(Y-2X+1)/4$, and the
retained square equation becomes $Y=X^2$. This converts
the old label to $51X^2+2504$. The scalar shift is $9792$.
After each row is made primitive with a fixed sign, their
sum of squares is expanded by exact integer arithmetic.

The independent checker performs the substitution in the
reverse direction for every unchanged row. It separately
checks the two new scale rows, the root-coordinate
normalizations, the sole uses of deleted coordinates,
all finite prime congruences and the exact coefficient
expansion. All 50 equations of the third detector recover
the corresponding second-detector equations under the inverse
coordinate permutation; every use of its 52 coordinates is
checked. The full rational normalizations displayed above
are also verified as exact polynomial identities.
It computes the sum of squares by triangular
convolution, independently of the producer's ordered-pair
expansion. A coefficient equal to one proves content one
for every integer parameter; a nonzero degree-four term
proves the degree is exactly four. A coefficient root
bound proves the stated support for $N\ge27464$, and direct
integer evaluation gives the stated binary size for the
explicit member.

\begin{center}
\begin{tabular}{lr}
\toprule
Quantity & Exact value\\
\midrule
Total rational coordinates &370\\
Quadratic equations &342\\
Polynomial degree &4\\
Terms before grouping powers of $A_N$ &5852\\
Coordinate monomials, $N\ge27464$ &5837\\
Content &1\\
Largest coefficient bits, $N=27464$ &1791357\\
\bottomrule
\end{tabular}
\end{center}

Reproduction requires only Python's standard library:
\begin{verbatim}
python3 build_quartic370.py
python3 check_quartic370.py
\end{verbatim}
The compressed exact source table is included. These finite
checks establish the algebraic specification and numerical
certificates. Propositions~\ref{prop:soundness} and
\ref{prop:complete}, together with the stated arithmetic
premises, establish the rational input projection.

\begin{thebibliography}{9}
\bibitem{Quartic392}
\emph{A forced prime condition and an explicit quartic in 392 variables},
8 November 2026, source arithmetic contract and exact circuit,
\url{https://github.com/JoseBrox/diophantine-z-q_improvements},
\nolinkurl{papers/quartic/quartic392_20261108_v8.tex}.
\bibitem{Quartic389}
\emph{A quartic in 389 variables from local nonvanishing},
8 November 2026, same repository,
\nolinkurl{papers/quartic/quartic389_20261108_v9.tex}.
\bibitem{Quartic317}
\emph{A family of quartics in 317 variables from aggregate trace heights},
8 November 2026, same repository,
\nolinkurl{papers/quartic_existence/quartic317_family_20261108_v9.tex}.
\bibitem{EffectiveSums}
\emph{Effective heights of sums of elliptic abscissas},
8 November 2026, companion paper,
\nolinkurl{effective_elliptic_sums_20261108_v10.tex}.
\end{thebibliography}
\end{document}
